\FloatBarrier
\section{Conclusion}

This paper began from the observation that the binding constraint in applied causal inference is not estimation but the graph: every tool in the Pearlian framework is conditional on a causal structure that, in most observational work, is asserted rather than tested. Its answer is a class and a procedure. Compliance-gated administrative datasets, defined by the three conditions of Section 2 (unit-level threshold selection on observables, endogeneity of the selection variable, atomic reporting), are a class for which that problem is tractable, because the document that creates the compliance obligation also fixes the structure it imposes on the data. The TDGP pipeline turns that property into a procedure: locate the gate in the missingness pattern (Stage 1), formulate the DAG from the gate, the shadow matrix, and the variable descriptions (Stage 2), and expose the hypothesis to falsification against the constitutive documents, which were held out from its construction (Stage 3).

On the Chicago Energy Benchmarking data the procedure ran end to end. Stage 3's sharpest test reconstructed the recorded outcome as the documented linear combination of the metered fuels to $R^2 = 0.9999$, leaving no role for a parent of the outcome beyond the mediators, and the surviving structure then licensed the ladder's upper rungs in turn. At the rung of intervention, the inverse-probability correction implemented the back-door adjustment the graph certified, and the 2$\times$3 factorial isolated what it contributes: on the corrected foundation, the structural Bayesian model matched the unconstrained gradient-boosted ensemble to within 0.4 percentage points of median absolute percentage error (21.4 against 21.0). The tie is the point, because of what it purchases: a floor-area elasticity of 0.979 with an interval that excludes proportional scaling, a decarbonisation trend of roughly 5.4 percent per year that survived every prior and weighting specification tested, a 22-fold emissions range across property types beyond what size, age, and reporting year explain, and uncertainty-carrying emission estimates for the 5,510 building-year records that never reported under the ordinance. At the rung of imagining, the same posterior answered counterfactual queries with outputs the point-prediction pipelines of Section 5 do not deliver: under trend continuation to 2030, 37.1 percent of records remain at high risk against a stylised compliance threshold, and reassigning every building's vintage to 2010 credibly raises predicted emissions for 90.2 percent of records, a sign the tested separation of trend from cohort makes interpretable rather than suspect.

The convergence of the two model families is a corollary, not a coincidence. \cit{breiman2001}{Breiman (2001)} framed the structural and algorithmic cultures as rivals; the factorial shows the rivalry dissolving under a specific and checkable condition, namely that the data-generating process is recoverable from institutional documentation and has survived a falsification test against it. When that condition holds, both families estimate the same population quantity and differ only in what they return: a number, or a mechanism with a number attached. When it does not hold, the algorithmic culture's caution stands, and nothing in this paper argues otherwise.

\subsection{Limitations}

\textbf{Selection ignorability is the residual the documents cannot reach.} Stage 3 tests which variables the gate keys on, and the data show that realised compliance departs from the documented rule: the propensity model classifies compliance at 95.1 percent accuracy, not 100 (Section 5.3). Wherever it departs, the correction rests on the standard assumption that no unobserved driver of compliance is correlated with the outcome given the upstream features (\cit{hernanrobins2020}{Hernán \& Robins, 2020}). The gate-fidelity and manipulation checks the pipeline affords (Sections 3.1, 5) audit this partially, but the unobserved component of selection that Section 5 carries as a limitation (management practices, voluntary retrofits, tenant behaviour) is exactly what no within-sample diagnostic can rule out.

\textbf{The correction stops at the overlap boundary.} Inverse-probability weighting recovers the population only where compliant units exist to be reweighted. Below the 50,000 square foot threshold there are none, so every statement about the never-observed region rides on the modularity of the tested equations and on the assumption that the learned upstream form is stable past the cutoff. Even within support, where the 5,510 non-reporting records sit, prediction outruns validation: the deployment inference of Section 6.4 was framed as consistency with the tested DGP, not as validation against ground truth, because for a never-reporting building no ground truth exists.

\textbf{The temporal trend is a bundle, and the projections assume its persistence.} $\beta_T$ mixes the decline in the grid emission factor with efficiency change in the stock, and any panel-wide operational shifts are absorbed into it as well, since occupancy and operational schedules are unobserved (Section 6.2). The policy comparisons of Section 7.1 inherit both qualifications: the pace that clears the city's 2040 target is a building-stock rate carried forward unchanged, not a forecast.

\textbf{The pipeline tests structure, not functional form.} Stages 1 to 3 settle which variables enter and how observation is gated; the linear world-layer form of the focal model was a downstream choice, and the factorial itself showed that form still matters when it is wrong: the causally-blind Bayesian model could not exploit its access to the mediators because a log-linear likelihood cannot represent their additive accounting identity (Section 5.6). The convergence result is therefore conditional on the chosen form being adequate for this outcome, as the predictive tie indicates it is here, and nothing in the pipeline guarantees adequacy elsewhere.

\textbf{Mechanism-level counterfactuals remain out of reach.} The energy mediators are compliance-gated, so the focal model operates on the reduced form from upstream features to outcome. Queries that act on the mediators, retrofit impact above all, are not identified by anything estimated here, a boundary Section 7.3 drew explicitly.

\textbf{The empirical evidence is one member of the class.} Section 2.4 argued the class is broad; the demonstration is a single city's programme over a single decade. Whether the pipeline ports as claimed, and whether parameters such as the near-unit elasticity generalise, are empirical questions this paper poses but cannot answer.

\subsection{Extensions}

The limitations order the research programme. The nearest extension is replication across the class: benchmarking ordinances of the same gate structure operate in other cities, differing in thresholds and stringency, so the pipeline can be re-run end to end and the structural parameters compared across programmes, turning the single-city caveat into a test of generalisation. Section 2.4 points further afield, to injury recordkeeping, toxics release reporting, and credit origination, where the constitutive documents differ in kind but the three conditions hold in the same form.

Two extensions strengthen the correction itself. Replacing the single propensity model with doubly robust estimation would let the analysis survive misspecification of either the selection or the outcome model; it would not discharge the ignorability assumption, which no estimator removes, but it would stop the correction depending on one specification being right. And modelling the mediator layer, which the data record for the (largely) compliant units, would extend the reach of the third rung: an inverse-probability-corrected upstream-to-mediator model, composed with the documented accounting identity, is exactly what a retrofit counterfactual requires, converting Section 7.3's stated boundary into a specification.

The deployment direction runs toward the regulator. The compliance-risk machinery of Section 7.2 applies unchanged to whatever threshold a performance standard specifies, and the posterior-predictive inventory of Section 6.4 replaces point-imputed emission gap-filling with estimates that state their uncertainty. A study built with a city's actual enforcement parameters would test whether the probabilistic classification earns its keep in practice.

The general claim survives the particulars. Wherever a written rule fixes which units are recorded, and that rule can be read, the DAG establishment problem admits a test, and a graph that has survived a test licenses what an asserted graph only borrows: adjustment with a warrant, extrapolation with a stated cost, and counterfactuals whose surprises are findings rather than artefacts (\cit{pearl2009}{Pearl, 2009}). The compliance gate, ordinarily a reason to distrust administrative data, is on this reading its most valuable feature, because the same institution that censors the data documents the process that generated it.

\subsection{Data and Code Availability}

The Chicago Energy Benchmarking dataset is publicly available from the City of Chicago Data Portal (\cit{citynd}{City of Chicago, n.d.}) and was retrieved through the portal's API; no data is redistributed with this paper. The analysis code implementing the TDGP pipeline, the 2$\times$3 factorial experiment, and the counterfactual queries is available from the author on request.

